\section{Hasil Ekstraksi Aspek Menggunakan BERTopic}

BERTopic digunakan untuk mengidentifikasi aspek atau topik yang paling sering dibahas pada ulasan Google Reviews. Hasil pemodelan menghasilkan 18 topik yang kemudian diberi label secara manual berdasarkan kata-kata dominan pada masing-masing topik.

\begin{table}[H]
\centering
\caption{Daftar Topik Hasil BERTopic}
\label{tab:hasil_bertopic}
\begin{tabular}{|c|p{9cm}|}
\hline
\textbf{ID Topik} & \textbf{Label Topik} \\
\hline
-1 & Lainnya / Noise \\
0 & Keindahan dan Spot Foto \\
1 & Wisata Religi (Makam Syaikhona Kholil) \\
2 & Wisata Alam (Bukit Kapur Jaddih) \\
3 & Wisata Religi dan Kegiatan Umum \\
4 & Lainnya / Noise \\
5 & Keluhan Pungli dan Pengemis \\
6 & Infrastruktur dan Akses Jalan \\
7 & Lanskap Tambang Kapur \\
8 & Lainnya / Noise \\
9 & Kondisi Alam dan Danau \\
10 & Lainnya / Noise \\
11 & Kebersihan dan Area Asri \\
12 & Keindahan Pemandangan \\
13 & Kepadatan Peziarah \\
14 & Lainnya / Noise \\
15 & Fasilitas (Parkir dan Toilet) \\
16 & Suasana Tenang dan Damai \\
17 & Penataan Lokasi \\
\hline
\end{tabular}
\end{table}

Berdasarkan hasil BERTopic, aspek yang paling banyak dibahas meliputi keindahan wisata, fasilitas, akses jalan, kondisi lingkungan, dan wisata religi. Selain itu, ditemukan pula beberapa topik yang tergolong \textit{noise} karena tidak memiliki pola pembahasan yang jelas. Hasil ekstraksi aspek ini selanjutnya digunakan dalam proses \textit{Aspect-Based Sentiment Analysis} (ABSA) bersama hasil klasifikasi sentimen dari IndoBERT.